\documentclass[12pt,a4paper]{article}
\usepackage[utf8]{inputenc}
\usepackage[T1]{fontenc}
\usepackage[spanish,es-noindentfirst]{babel}
\usepackage[hidelinks]{hyperref}
\usepackage{booktabs}
\usepackage{array}
\usepackage{graphicx}
\usepackage{caption}
\usepackage{amsmath}
\usepackage{amssymb}
\usepackage{geometry}
\geometry{margin=2.5cm}
\renewcommand{\familydefault}{\sfdefault}
\setlength{\parskip}{0.5em}
\setlength{\parindent}{0pt}

\title{Marco de Criterios, Métricas y Modelo de Decisión para la Selección Objetiva de CNI en Kubernetes}
\author{Holman Alba \and Duwan Sierra}
\date{Mayo 2026}

\begin{document}

\maketitle

\begin{abstract}
Este documento expone el marco técnico del \textbf{Objetivo Específico 3} del trabajo de grado: definir criterios, métricas y factores de ponderación para la comparación objetiva de plugins CNI en Kubernetes.
Se presentan la metodología MCDA, los umbrales de rendimiento, la escala de puntuación y un primer análisis con datos reales de un clúster K3s en DigitalOcean.
\end{abstract}

\section{Marco Conceptual}

El modelo propuesto transforma los datos brutos obtenidos del benchmarking automatizado en una \textbf{decisión arquitectónica objetiva}. El proceso se divide en cuatro fases:

\begin{enumerate}
  \item Datos crudos (JSON/CSV)
  \item Limpieza IQR mediante \texttt{procesador.js}
  \item Métricas consolidadas y umbrales normalizados
  \item Scoring y ponderación por perfil
\end{enumerate}

El resultado es un puntaje final para cada CNI que permite comparar alternativas sin depender de juicios subjetivos.

\subsection{Formalización del modelo MCDA}

La metodología se asienta en el análisis multicriterio (MCDA), que transforma cada métrica en un valor numérico y lo combina mediante factores de peso.
La fórmula general es:

\begin{equation}
P_{CNI} = \sum_{i=1}^{n} V_i \times W_i
\end{equation}

Donde:
\begin{itemize}
  \item $P_{CNI}$ es el puntaje total para un plugin de red.
  \item $V_i$ es la puntuación de la métrica $i$ según cumplimiento de umbrales.
  \item $W_i$ es el peso de la métrica $i$ en el caso de uso.
  \item $n$ es el número total de métricas consideradas.
\end{itemize}

\section{Criterios y Métricas}

El modelo organiza los indicadores en tres criterios principales:

\begin{itemize}
  \item \textbf{C1 -- Rendimiento de red}
  \item \textbf{C2 -- Eficiencia de recursos}
  \item \textbf{C3 -- Impacto de seguridad}
\end{itemize}

\begin{table}[htbp]
\centering
\caption{Criterios y métricas del modelo de evaluación.
}\label{tab:criterios-metricas}
\begin{tabular}{p{2.3cm}p{5.5cm}p{6cm}}
\toprule
\textbf{Criterio} & \textbf{Métricas} & \textbf{Fuente de datos} \\
\midrule
C1 & Throughput TCP, latencia media, latencia máxima, jitter MDEV, retransmisiones TCP & \texttt{iperf3} JSON por CronJob \\
C2 & CPU del nodo, RAM por nodo & Exportadores de Prometheus / Grafana, CSV \\
C3 & Overhead de latencia con NetworkPolicy, overhead de throughput con NetworkPolicy & Cálculo automático en \texttt{procesador.js} \\
\bottomrule
\end{tabular}
\end{table}

\section{Escala de Puntuación}

Cada métrica se traduce a una escala de tres valores:

\begin{itemize}
  \item \textbf{5 puntos}: Excelente, cumple los requisitos más estrictos.
  \item \textbf{3 puntos}: Aceptable, opera dentro de límites funcionales.
  \item \textbf{1 punto}: Deficiente, no cumple el umbral del caso de uso.
\end{itemize}

Esta traducción convierte mediciones heterogéneas en valores comparables matemáticamente.

\section{Umbrales de Rendimiento}

Los umbrales se adaptan a tres casos de uso arquitectónico, diferenciando aplicaciones con requisitos de latencia, eficiencia o seguridad.

\subsection{Caso de Uso 1: URLLC / Automatización Industrial}

\begin{table}[htbp]
\centering
\caption{Umbrales para URLLC / automatización industrial.}
\label{tab:umbrales-urllc}
\begin{tabular}{p{3.5cm}p{3.5cm}p{3.5cm}p{3cm}}
\toprule
\textbf{Métrica} & \textbf{Excelente (5)} & \textbf{Aceptable (3)} & \textbf{Deficiente (1)} \\
\midrule
Latencia promedio (ms) & < 1 & 1--10 & > 10 \\
Jitter MDEV (ms) & < 0.1 & 0.1--1 & > 1 \\
Throughput TCP (Mbps) & > 1000 & 100--1000 & < 100 \\
CPU overhead del CNI (\%) & < 1 & 1--5 & > 5 \\
RAM por nodo (MB) & < 50 & 50--200 & > 200 \\
Overhead NP latencia (\%) & < 5 & 5--15 & > 15 \\
\bottomrule
\end{tabular}
\end{table}

\subsection{Caso de Uso 2: Edge Computing / IoT}

\begin{table}[htbp]
\centering
\caption{Umbrales para Edge Computing / IoT.}
\label{tab:umbrales-edge}
\begin{tabular}{p{3.5cm}p{3.5cm}p{3.5cm}p{3cm}}
\toprule
\textbf{Métrica} & \textbf{Excelente (5)} & \textbf{Aceptable (3)} & \textbf{Deficiente (1)} \\
\midrule
Latencia promedio (ms) & < 10 & 10--50 & > 50 \\
Jitter MDEV (ms) & < 1 & 1--5 & > 5 \\
Throughput (% de nominal) & > 80\% & 50--80\% & < 50\% \\
CPU overhead del CNI (\%) & < 5 & 5--15 & > 15 \\
RAM por nodo (MB) & < 50 & 50--150 & > 150 \\
Overhead NP latencia (\%) & < 10 & 10--20 & > 20 \\
\bottomrule
\end{tabular}
\end{table}

\subsection{Caso de Uso 3: Microservicios Transaccionales / Zero Trust}

\begin{table}[htbp]
\centering
\caption{Umbrales para microservicios transaccionales / Zero Trust.}
\label{tab:umbrales-zero}
\begin{tabular}{p{3.5cm}p{3.5cm}p{3.5cm}p{3cm}}
\toprule
\textbf{Métrica} & \textbf{Excelente (5)} & \textbf{Aceptable (3)} & \textbf{Deficiente (1)} \\
\midrule
Latencia promedio (ms) & < 50 & 50--200 & > 200 \\
Jitter MDEV (ms) & < 5 & 5--20 & > 20 \\
Throughput (% de nominal) & > 70\% & 50--70\% & < 50\% \\
CPU overhead del CNI (\%) & < 8 & 8--15 & > 15 \\
RAM por nodo (MB) & < 200 & 200--400 & > 400 \\
Overhead NP latencia (\%) & < 5 & 5--15 & > 15 \\
Overhead NP throughput (\%) & < 5 & 5--15 & > 15 \\
\bottomrule
\end{tabular}
\end{table}

\section{Datos Reales del Testbed}

En este documento se utilizan datos de un clúster K3s en DigitalOcean, con 5 corridas de benchmarking de 300 segundos por CNI y limpieza IQR automática.

\begin{table}[htbp]
\centering
\caption{Resultados reales de benchmark por CNI.}
\label{tab:datos-reales}
\begin{tabular}{p{4.6cm}p{2cm}p{2cm}p{2cm}p{2cm}}
\toprule
\textbf{Métrica} & \textbf{Flannel} & \textbf{Calico} & \textbf{Cilium} & \textbf{Antrea} \\
\midrule
Throughput TCP (Mbps) & 1,351 & 1,816 & 931 & 1,291 \\
Retransmisiones TCP & 6,526 & 86,651 & 12,975 & 21,863 \\
Latencia promedio (ms) & 22.7 & 20.5 & 29.5 & 41.3 \\
Latencia máxima (ms) & 52 & 46 & 99 & 112 \\
Jitter MDEV (ms) & 8.66 & 6.00 & 16.00 & 19.36 \\
CPU del CNI (\%) & Pendiente & Pendiente & Pendiente & Pendiente \\
RAM por nodo (MB) & Pendiente & Pendiente & Pendiente & Pendiente \\
Overhead NP latencia (\%) & N/A & Pendiente & Pendiente & Pendiente \\
Overhead NP throughput (\%) & N/A & Pendiente & Pendiente & Pendiente \\
\bottomrule
\end{tabular}
\end{table}

\section{Matrices de Scoring}

Los umbrales anteriores se aplican a los datos reales para obtener puntuaciones que permiten comparar los CNIs.

\subsection{Ejemplo de scoring parcial para Edge/IoT}

El perfil Edge/IoT utiliza mayor peso en eficiencia de recursos y una escala de latencia moderada.

\begin{table}[htbp]
\centering
\caption{Puntaje parcial de ejemplo para Edge/IoT (solo métricas de red disponibles).}
\label{tab:scoring-edge}
\begin{tabular}{p{3cm}p{2cm}p{2cm}p{2cm}p{2cm}}
\toprule
\textbf{CNI} & \textbf{Latencia} & \textbf{Jitter} & \textbf{Throughput} & \textbf{Subtotal} \\
\midrule
Flannel & 5 & 1 & 3 & 9 \\
Calico & 5 & 1 & 5 & 11 \\
Cilium & 5 & 1 & 1 & 7 \\
Antrea & 5 & 1 & 3 & 9 \\
\bottomrule
\end{tabular}
\end{table}

\section{Factores de Ponderación}

Los pesos reflejan la prioridad técnica de cada caso de uso.

\begin{table}[htbp]
\centering
\caption{Pesos de métricas para cada perfil arquitectónico.}
\label{tab:pesos-perfiles}
\begin{tabular}{p{3.5cm}p{2.5cm}p{2.5cm}p{2.5cm}}
\toprule
\textbf{Métrica} & \textbf{URLLC} & \textbf{Edge/IoT} & \textbf{Zero Trust} \\
\midrule
Latencia promedio & 0.30 & 0.20 & 0.15 \\
Jitter / pérdida & 0.20 & 0.10 & 0.20 \\
Throughput & 0.20 & 0.10 & 0.15 \\
CPU/RAM & 0.20 & 0.35 & 0.15 \\
Overhead NP & 0.10 & 0.10 & 0.40 \\
\bottomrule
\end{tabular}
\end{table}

\section{Puntaje Parcial con Datos Disponibles}

Con los datos actuales, solo es posible calcular una aproximación parcial. Estos puntajes muestran la dirección técnica de la comparación, pero dependen de la incorporación futura de CPU, RAM y métricas de Network Policies.

\section{Actualización Automática del Modelo}

El pipeline automatizado está diseñado para actualizarse cuando se agregan nuevos datos:

\begin{enumerate}
  \item CSV de uso de recursos en \texttt{resource_usage_nodes/}
  \item JSON de overhead con Network Policy en \texttt{with_network_policy/}
  \item Ejecución de \texttt{procesador.js} para recalcular los puntajes
  \item Publicación automática del conjunto de datos actualizado en GitHub Pages
\end{enumerate}

\section{Referencias}

\begin{itemize}
  \item ITU-R M.2410-0. Requisitos mínimos para IMT-2020.
  \item 3GPP TS 22.104. Service Requirements for Cyber-Physical Control.
  \item 3GPP TS 22.261. Service requirements for the 5G system.
  \item NIST SP 800-207. Zero Trust Architecture.
  \item Karampelas et al. (2024). Performance Evaluation of Kubernetes Networking Approaches.
  \item ACM 10.1145/3479645.3479700. Performance Analysis of CNI Plugins.
\end{itemize}

\end{document}
